\documentclass[]{article}

%opening
\title{Brusselator Reaction Finite Difference Derivation and Documentation}
\author{Griffin Layhew}

\begin{document}

\maketitle

\begin{abstract}

\end{abstract}

\section{Problem Definition}
The Brusselator is a model for an autocatalytic reaction, which can be written as a reaction-diffusion equation for two chemical species, U and V, which each have different diffusion rates. The system is defined as follows:

\begin{equation}
	\frac{\partial U(\mathbf{x}, t)}{\partial t} = D_U \nabla^2 U(\mathbf{x}, t) + A + U(\mathbf{x}, t)^2 V(\mathbf{x}, t) - (B+1) U(\mathbf{x}, t)
\end{equation}

\begin{equation}
	\frac{\partial V(\mathbf{x}, t)}{\partial t} = D_V \nabla^2 V(\mathbf{x}, t) + BU(\mathbf{x}, t) - U(\mathbf{x}, t)^2 V(\mathbf{x}, t)
\end{equation}

\begin{equation}
	U(\mathbf{x}, t) = A + \sigma
\end{equation}

\begin{equation}
	V(\mathbf{x}, t) = B + \sigma
\end{equation}

\begin{equation}
	\left. \frac{\partial U(\mathbf{x}, t)}{\partial t} \right|\Gamma = 0 
\end{equation}

\begin{equation}
	\left. \frac{\partial V(\mathbf{x}, t)}{\partial t} \right|\Gamma = 0 
\end{equation}



\section{Spatial Discretization (Finite Difference Method)}
The Brusselator system defined earlier can be decomposed using the finite difference method. Using this method...

\begin{equation}
	\frac{U^{n+1}_{j,k} - U^n_{j,k}}{\Delta t} = D_U \left( \frac{U^{n}_{j+1,k} - 2 U^{n}_{j,k} + U^{n}_{j-1,k}}{\Delta x^2} + \frac{U^{n}_{j,k+1} - 2 U^{n}_{j,k} + U^{n}_{j,k-1}}{\Delta y^2}\right) + A + {U^{n}_{j,k}}^2 V^{n}_{j,k} - (B+1) U^{n}_{j,k}
\end{equation}

\begin{equation}
	\frac{V^{n+1}_{j,k} - V^n_{j,k}}{\Delta t} = D_V \left( \frac{V^{n}_{j+1,k} - 2 V^{n}_{j,k} + V^{n}_{j-1,k}}{\Delta x^2} + \frac{V^{n}_{j,k+1} - 2 V^{n}_{j,k} + V^{n}_{j,k-1}}{\Delta y^2}\right) + B U^{n}_{j,k} - {U^{n}_{j,k}}^2 V^{n}_{j,k}
\end{equation}

\section{Temporal Approximation (Explicit RK4 Method)}
It is helpful to defined the spatial derivative operators for both U and V as follows:
\begin{equation}
	\mathcal{L}(U^{n}_{j,k}) = D_U \left( \frac{U^{n}_{j+1,k} - 2 U^{n}_{j,k} + U^{n}_{j-1,k}}{\Delta x^2} + \frac{U^{n}_{j,k+1} - 2 U^{n}_{j,k} + U^{n}_{j,k-1}}{\Delta y^2}\right) + A + {U^{n}_{j,k}}^2 V^{n}_{j,k} - (B+1) U^{n}_{j,k}
\end{equation}

\begin{equation}
	\mathcal{L}(V^{n}_{j,k}) = D_V \left( \frac{V^{n}_{j+1,k} - 2 V^{n}_{j,k} + V^{n}_{j-1,k}}{\Delta x^2} + \frac{V^{n}_{j,k+1} - 2 V^{n}_{j,k} + V^{n}_{j,k-1}}{\Delta y^2}\right) + B U^{n}_{j,k} - {U^{n}_{j,k}}^2 V^{n}_{j,k}
\end{equation}

Then, the RK4 method for is defined as follows:
% Step 1 of Runge Kutta Method for Coupled Systems
\begin{center}
	\textbf{Step 1}
\end{center}

\begin{equation}
	U^{1}_{j,k} = U^{n}_{j,k} + \frac{\Delta t}{4} \mathcal{L}(U^{n}_{j,k})
\end{equation}
\begin{equation}
	V^{1}_{j,k} = V^{n}_{j,k} + \frac{\Delta t}{4} \mathcal{L}(V^{n}_{j,k})
\end{equation}

% Step 2 of Runge Kutta Method for Coupled Systems
\begin{center}
	\textbf{Step 2}
\end{center}

\begin{equation}
	U^{2}_{j,k} = U^{n}_{j,k} + \frac{\Delta t}{3} \mathcal{L}(U^{1}_{j,k})
\end{equation}
\begin{equation}
	V^{2}_{j,k} = V^{n}_{j,k} + \frac{\Delta t}{3} \mathcal{L}(V^{1}_{j,k})
\end{equation}

% Step 3 of Runge Kutta Method for Coupled Systems
\begin{center}
	\textbf{Step 3}
\end{center}

\begin{equation}
	U^{3}_{j,k} = U^{n}_{j,k} + \frac{\Delta t}{2} \mathcal{L}(U^{2}_{j,k})
\end{equation}
\begin{equation}
	V^{3}_{j,k} = V^{n}_{j,k} + \frac{\Delta t}{2} \mathcal{L}(V^{2}_{j,k})
\end{equation}

% Step 4 of Runge Kutta Method for Coupled Systems
\begin{center}
	\textbf{Step 4}
\end{center}

\begin{equation}
	U^{n+1}_{j,k} = U^{n}_{j,k} + \Delta t \mathcal{L}(U^{3}_{j,k})
\end{equation}
\begin{equation}
	V^{n+1}_{j,k} = V^{n}_{j,k} + \Delta t \mathcal{L}(V^{3}_{j,k})
\end{equation}



\end{document}
